\documentclass[12pt]{extarticle}

\usepackage[T1]{fontenc}
\usepackage{polski}
\usepackage[utf8x]{inputenc}
\usepackage[polish]{babel}
\usepackage{url}
\usepackage{afterpage}

\usepackage[margin=1.0in]{geometry}
\usepackage{color}
\usepackage{graphicx}
\usepackage{float}
\usepackage{indentfirst}
\setlength\parindent{1cm}

\usepackage{listings}
\lstset{
	language=Java,
	basicstyle=\ttfamily,
	numbers=left,
 	numberstyle=\tiny,
	frame=tb,
	tabsize=4,
	columns=fixed,
	showstringspaces=false,
	showtabs=false,
	keepspaces,
	commentstyle=\color{red},
	keywordstyle=\color{blue}
}

\begin{document}
\begin{titlepage}
    \begin{center}
        \includegraphics[width=4cm]{polsl.png}\\[1cm]
        \textsc{\LARGE{Politechnika Śląska}}\\[0.5cm]
        \textsc{\LARGE{Wydział Automatyki, Elektroniki i~Informatyki}}\\[0.5cm]
        \textsc{\LARGE{Kierunek Informatyka}}\\[2.5cm]
        \LARGE{Praca magisterska}\\[1cm]
        \begingroup
            \fontsize{14pt}{17pt}\selectfont
            Analiza porównawcza przetwarzania strumieniowego w Java 8 \\ i realizacji zapytań SQL w pamięci
        \endgroup
    \end{center}
    \vspace{2cm}
    \begingroup
        \fontsize{14pt}{17pt}\selectfont
        \textbf{Autor:} Aleksander Grzybowski\\
        \textbf{Kierujący pracą:} dr inż. Ewa Płuciennik\\
    \endgroup

    \vspace{1.0cm}
    \begingroup
        \fontsize{12pt}{14pt}\selectfont
        \begin{center}
        Gliwice, kwiecień 2017
        \end{center}
    \endgroup
\end{titlepage}

\clearpage

\tableofcontents

\newpage

\section{Metody dostępu do danych}

\subsection{Bazy danych SQL w pamięci (in-memory)}

    Bazy danych SQL są powszechnie stosowanymi rozwiązaniami do gromadzenia i odpytywania danych. Struktura danych ściśle określona przez definicje tabel, ochrona integralności poprzez ograniczenia i więzy referencyjne oraz wygodny język SQL zachęcają do ich stosowania dla zarówno małych, jak i dużych zbiorów danych. Mimo pojawienia się nowych technologii, takich jak bazy dokumentowe typu MongoDB czy bazy grafowe (Neo4j) tradycyjne bazy relacyjne są wciąż najczęściej stosowanym rozwiązaniem. Bazy danych działające całkowicie w pamięci RAM wykorzystywane są w dużych hurtowniach danych, gdzie szybkość ma pierwszorzędne znaczenie, a niemożliwe jest użycie innego rozwiązania, np. cache. 

    Jako iż porównanie baz danych i strumieni w Java 8 powinno być miarodajne, konieczne jest wybranie do badań baz in-memory działających w środowisku Java. Dwoma najpopularniejszymi rozwiązaniami są H2 oraz HSQLDB. Są one wykorzystywane głównie w celach deweloperskich, jako ulotne środowiska do testowania programów wymagających bazy danych SQL. Dzięki działaniu całkowicie w pamięci, nie jest konieczne żadne zarządzanie, a zawartość bazy jest tracona po zamknięciu programu. Sposób korzystania z tych baz jest identyczny, jak przy korzystaniu z tradycyjnych baz dyskowych, a interfejs programistyczny JDBC nie różni się w żaden sposób. Możliwe jest więc tez porównanie wydajności baz dyskowych z bazami in-memory, które także zostanie wykonane.
    TODO coś więcej o SQL z przykładami


\subsection{Strumienie w Java 8}

    Stream API jest wprowadzonym w wersji 8 Javy komponentem biblioteki standardowej, który pozwala na tworzenie tzw. strumieni. Strumień jest to abstrakcyjna reprezentacja ciągu obiektów, które pochodzą z jakiegoś źródła, zwykle kolekcji, ale mogą być też pobierane bądź tworzone na bieżąco. Na strumieniach możliwe jest wykonywanie prostych i złożonych operacji filtrowania, sortowania, mapowania i grupowania, jak i innych, zdefiniowanych przez użytkownika operacji. Należy mieć na uwadze, że wszystkie wykonane w ciągu operacje wykonywane są 'leniwie', tzn. żadna obróbka danych nie zachodzi przed wykonaniem tzw. operacji terminalnej (zwykle zebranie wyników), która łańcuchowo powoduje wykonanie się wszystkich operacji pośrednich.

    Cechą szczególną strumieni jest możliwość ich łatwego zrównoleglania, udostępniona programistom w postaci metody \texttt{Stream.parallelStream}. Wielowątkowe przetwarzanie danych jest dość trudnym zagadnieniem w Javie, wymagana jest duża wiedza na temat działania maszyny wirtualnej i szczególna ostrożność. Jako iż strumienie definiowane są zwykle przez funkcje nie zmieniające stanu, dodanie równoległości jest znacznie prostsze. Jednym z celów niniejszej pracy jest zbadanie, jaki wpływ na wydajność mogą mieć równoległe strumienie.

    Poniższy listing prezentuje przykładowe użycie strumieni. Dostępna jest kolekcja \texttt{Collection<Customer> customers}, która zawiera obiekty reprezentujące klientów. Każdy klient posiada pola danych informujące między innymi o kraju jego pochodzenia oraz stanie konta (jest to część modelu danych TPC, który zostanie omówiony szczegółowiej w rozdziale ...). Metoda \texttt{stream} tworzy strumień obiektów typu \texttt{Customer}, na którym dokonuje się operacji. Metoda \texttt{Stream.filter} przyjmuje obiekt funkcyjny typu \texttt{Predicate}, który określa, czy dany element powinien znaleźć się w wynikowym strumieniu. Warto zaznaczyć, że możliwe jest wykorzystanie dowolnego kodu języka Java - nie ma tutaj ograniczeń na konkretny zbiór funkcji, jak to jest w języku SQL. Metoda \texttt{sorted} zwraca nowy strumień, posortowany według podanego klucza - w domyślnej postaci porównywanie są bezpośrednio obiekty znajdujące się w strumieniu, tutaj wykorzystany został wariant metody sortującej z funkcją wyciągającą klucz sortowania z obiektu (tutaj - stan konta klienta). Na koniec wyniki zbierane są do listy typu \texttt{java.util.List}. 

\begin{lstlisting}[label=testclass, caption=Przykładowe wykorzystanie Stream API]

List<Customer> canadiansSortedByBalance = customers.stream()
        .filter(customer -> Objects.equals(
                customer.nation.name,
                "CANADA"
        ))
        .sorted(Comparator.comparing(
                customer -> customer.acctbal
        ))
        .collect(toList());

\end{lstlisting}

    Bardziej złożony przykład zaprezentowany jest na następnym listingu. Kolekcja zamówień, na której opiera się zapytanie, składa się z obiektów typu \texttt{Order} posiadająch status oraz przyporządkowanego do zamówienia klienta. Podobnie jak w poprzednim przykładzie, w pierwszym kroku przeprowadza się filtrowanie jedynie zamówień o wyższym priorytecie. Następnie wykorzystana jest metoda \texttt{collect}, której zadaniem jest zebranie i przetworzenie wyników w ostateczną postać wynikową. Wykorzystano tutaj dwa zagnieżdżone kolektory, z których pierwszy dokonuje grupowania zamówień na podstawie przypisanego klienta, a drugi kolektor przejmuje pogrupowane zamówienia i dla każdej grupy określa ich wielkość, zliczając wszystkie obiekty w grupie. W ten sposób wyznaczone zostają liczby ważnych zamówień dla poszczególnych klientów.

\begin{lstlisting}[label=testclass, caption=Zaawansowane wykorzystanie Stream API]

List<String> high = asList("1-URGENT", "2-HIGH");
Map<Customer, Long> orderCountsByCustomer = orders.stream()
        .filter(o -> high.contains(o.orderPriority))
        .collect(
                Collectors.groupingBy(order -> order.customer,
                        Collectors.counting()
                )
        );
        
        

\end{lstlisting}

\subsection{Różnice pomiędzy metodami}

    Pomiędzy przetwarzaniem strumieniowym, a bazami SQL działającymi w pamięci występuje bardzo dużo różnic, pomimo pozornie podobnej zasady działania. Celem obu podejść jest uzyskanie wygodnego i względnie szybkiego wykonywania zapytań, natomiast różnice w implementacji znacząco wpływają na ten cel.

    Głównym blokiem budującym w modelu relacyjnej bazy danych jest relacja (tabela) posiadająca kolumny i rekordy, a integralność danych zapewniona jest w pewnym stopniu przez ograniczenia, klucze i więzy referencyjne. Język SQL pozwala dzięki mechanizmowi JOIN na łatwe łączenie powiązanych ze sobą tabel i uzyskiwanie połączonych wyników. Jeśli baza danych została odpowienio zaprojektowana, uwzględniając wymaganą normalizację danych, to łączenie tabel na podstawie kluczy obcych jest naturalne dla programisty i proste w zrozumieniu. W przetwarzaniu strumieniowym pojęcie tabeli nie istnieje, ponieważ kontenerami danych są zwykłe kolekcje języka Java (listy, zbiory i mapy). O ile istnieją biblioteki bazodanowe typu jOOQ, które pozwalają na konwersję wyniku zapytania JDBC na strumień, to tak naprawdę operacje strumieniowe wykonują się zawsze na kolekcji w pamięci maszyny wirtualnej. W związku, aby operacje strumieniowe były wygodne dla programisty, model danych powinien uwzględniać powiązania między poszczególnymi obiektami w typowy dla języka Java sposób, poprzez referencje do innych obiektów (kompozycja). Łączenie kolekcji, choć możliwe, komplikuje znacznie zapytania i zmniejsza ich ogólną szybkość.

    Bazy in-memory wykorzystują wewnątrz struktury danych, które są ścisłym szczegółem implementacyjnym, niewidocznym dla programisty i ukrytym przed zmianami. Jest to w większości przypadków dobre rozwiązanie, pozwalające na odseparowanie wewnętrznej reprezentacji danych od metody dostępu do nich (przykładowo, dane w bazie H2 reprezentowane są przez obiekty klas \texttt{org.h2.mvstore.Page} i \texttt{org.h2.value.*}). Jednakże przy takim podejściu programista nie ma możliwości optymalizacji (o ile to możliwe) struktur danych i musi polegać na wykorzystanym w silniku bazy optymalizatorze zapytań. W przypadku strumieni wewnętrzne procesy zachodzące przy przetwarzaniu danych są ukryte, natomiast możliwe jest w każdym przypadku przepisanie części procesu na czystą Javę. Strumienie są bardzo wydajne, natomiast tak jak każda abstrakcja nad danymi wprowadzają pewnien niewielki narzut. W przypadku, gdy szybkość jest parametrem krytycznym, możliwe jest zastąpienie operacji klasy \texttt{Stream} zwykłymi konstrukcjami języka (pętle, warunki). 

    Przy wykonywaniu zapytania SQL z poziomu API JDBC nie istnieje żadne sprawdzanie poprawniości zapytania oraz jego składni w czasie kompilacji. Zapytanie wysyłane jest jako zwykły ciąg znaków do serwera bazodanowego, który następnie odczytuje zapytanie i przystępuje do jego interpretacji i wykonania. W związku z tym niemożliwe jest zapewnienie, czy dana tabela istnieje, albo czy dana funkcja SQL istnieje w danym oprogramowaniu bazy danych. Istnieją co prawda rozwiązania typu ORM, które chronią programistę przed niewłaściwym wykorzystaniem obiektu mapowanego na tabelę, albo edytory programistyczne z integracjami z serwerami baz, które umożliwiają przeskanowanie istniejących symboli (nazw tabel, kolumn itp.) i skorygowanie programisty, natomiast są to rozwiązania zewnętrzne, działające poza językiem Java. Przetwarzanie strumieniowe odbywa się w całości w języku Java, który jest silnie typowany, przez co niewłaściwe wykorzystanie konstrukcji języka zostanie automatycznie dostrzeżone przez kompilator. Przykładowo, wynikiem grupowania zbioru obiektów typu \texttt{T} na podstawie ich pola o typie \texttt{U} będzie zawsze \texttt{Map<T, List<U>>}. Dzięki temu typowe błędy przy odczytywaniu wyników z zapytania SQL przy użyciu obiektu \texttt{ResultSet} (niewłaściwy typ kolumny, niestniejące pole) nie występują.

TODO funkcje SQL (vendor-specific) vs dowolny kod w Javie debugowalny
    

\section{Badanie wydajności}

\subsection{Kryteria oceny}

    W niniejszej pracy pod pojęciem 'wydajność' rozumie się efektywny, uśredniony czas wykonania się danego zapytania i uzyskania wyników dostępnych w pamięci programu. Kryterium czasowe zostało przyjęte ze względu na łatwość jego pomiaru i bezpośrednie jego przełożenie na ogólną wydajność systemu wykorzystującego dane zapytanie. Niemniej jednak, w celach porównawczych, wykonane zostały również pomiary ilości zajętej pamięci, w celu zbadania wpływu struktur danych baz in-memory na zużycie pamięci. Jednym z problemów może być działanie odśmiecacza (garbage collector), który może być uruchamiany przez maszynę wirtualną Java w różnych fazach wykonywania badań, tak więc konieczna jest szczególna ostrożność przy interpretowaniu wyników.
    Jedną z trudności, które napotkano przy wykonywaniu eksperymentów, jest znaczny wpływ zastosowanej platformy sprzętowej na uzyskane wyniki. Oczywiście wydajność procesora będzie różnicowała ilość operacji na sekundę, natomiast konieczne jest sprawdzenie, czy relatywne porównanie wyników na różnych platformach jest jest takie same. Jeśli tak, to możliwe jest wyciągnięcie bardziej ogólnych wniosków z uzyskanych wyników, w przeciwnym przypadku możliwe jest, że dobór platformy sprzętowej jest istotny.
    Przy porównywaniu dwóch zupełnie różnych metod dostępu do danych, konieczne jest ustalenie postaci wynikowej pojedynczego eksperymentu. W przypadku języka SQL i metod JDBC wynik zapytania do bazy danych reprezentowany jest przez obiekt \texttt{ResultSet}, który udostępnia metody do poruszania się po wynikowych wierszach i wyciągania pojedynczych kolumn. W przypadku Stream API, rezultat zapytania dostępny jest od razu w formie obiektów standardowych kolekcji języka Java, takich jak \texttt{List<T>} czy \texttt{Map<K,V>}. W związku z tym, aby zapewnić sprawiedliwe porównanie obu metod, po wykonaniu zapytania SQL konieczne jest rozpakowanie obiektu \texttt{ResultSet} i przeniesienie danych do kolekcji Javy, aby uzyskać wynik w tej samej postaci.

\section{Środowisko badawcze}

\subsection{Budowa środowiska badawczego}

    Aby zebrane wyniki były jak najbardziej miarodajne, konieczne jest utworzenie solidnego środowiska testowego, w którym możliwe będzie wykonanie wszystkich potrzebnych badań. Jako iż elementem badań wydajności będą fragmenty kodu, pożądane jest, by pomiary czasowe odbywały się w miarę możliwości w sposób zautomatyzowany, generując pod koniec szczegółowy raport.

    JMH (Java Microbenchmark Harness) jest najpopularniejszym narzędziem służącym do tzw. microbenchmarkingu, czyli badania wydajności czasowej krótkich fragmentów kodu. Typowe metody pomiaru czasu w kodzie maszyny wirtualnej Javy nie są dokładne ani odporne na błędy, przykładowo często stosowane wywołanie \texttt{System.currentTimeMillis} nie jest odporne na zmiany czasu systemowego w trakcie wykonywania badania i ma różną dokładność na różnych systemach operacyjnych, natomiast \texttt{System.nanoTime()} nie jest bezpieczna wątkowo na poziomie procesora, przez co wykonywanie się zadań testowych na kilku rdzeniach procesora może powodować mylące wyniki. Co więcej, ręczne wykonywanie pomiarów jest dość żmudne i wymagałoby stworzenia własnego narzędzia. Zdecydowano się więc na wykorzystanie narzędzia gotowego. Przy użyciu JMH możliwe jest proste definiowanie metod, których wydajność (wykonania na sekundę) będzie badana, możliwe jest określenie ilości iteracji, ustalanie zmienych parametrów (np. ilość wierszy danych testowych). Raport wynikowy zawiera szczegółowe informacje na temat ilości wykonań metody, wariancje czasów oraz inne statystyki, na podstawie których możliwe jest wyciągnięcie potrzebnych wniosków.

    Do budowania projektu wykorzystano Gradle, które jest popularnym i używanym w zastosowaniach produkcyjnych narzędziem budowania projektów w Javie. Służy ono do uruchamiania kompilacji, uruchamiania testów jednostkowych, jak i wykonywania testów wydajnościowych dzięki integracji z JMH, które dostępne jest w ekosystemie Gradle jako gotowy plugin z możliwością dowolnej konfiguracji.

\subsection{Schemat klasy testowej (microbenchmark)}

\begin{lstlisting}[label=testclass, caption=Przykładowa klasa JMH]

@State(Scope.Thread)
public class SampleMicrobenchmark {
    
    private Connection connection;
    private List<Integer> numbers = new ArrayList<>();
    
    @Param({"1000", "10000", "100000", "1000000"})
    public int numberCount;
    
    @Setup
    public void setup() throws Exception {
        connection = newDatabase();
        connection.createStatement().execute("CREATE TABLE numbers (val INT)");
        Random random = new Random();
        
        for (int i = 0; i < numberCount; ++i) {
            int randomNumber = random.nextInt(1000);
            numbers.add(randomNumber);
            String query = "insert into numbers values (" + randomNumber + ")";
            connection.createStatement(query).execute();
        }
    }
    
    @Benchmark
    public int loop_over_list() throws Exception {
        int sum = 0;
        for (int i : numbers) {
            sum += i;
        }
        return sum;
    }
    
    @Benchmark
    public int sql() throws Exception {
        ResultSet resultSet = connection.createStatement()
                .executeQuery("SELECT sum(val) FROM numbers");
        resultSet.next();
        return resultSet.getInt(1);
    }
    
    @Benchmark
    public int stream() throws Exception {
        return numbers.stream().mapToInt(i -> i).sum();
    }
    
    @Benchmark
    public int parallelStream() throws Exception {
        return numbers.parallelStream().mapToInt(i -> i).sum();
    }
}


\end{lstlisting}

    Każda zbiór metod testujących wydajność danego rozwiązania zawarty jest w jednej klasie testowej. Metoda \texttt{setup} przygotowuje potrzebne dane testowe - generując je w czasie wykonania programu w przypadku danych losowych, bądź wczytuje z pliku w przypadku wykorzystania benchmarka TPC. Dane testowe ładowane są do bazy SQL oraz do zwykłych kolekcji języka Java, głównie list. Pozostałe metody, oznaczone adnotacją \texttt{@Benchmark} są metodami przeznaczonymi do profilowania. Ważnym elementem jest zwrócenie wyniku z takiej metody, ponieważ w przeciwnym razie maszyna wirtualna Javy, widząc brak zwrócenia wyniku, może dokonać optymalizacji i usunąć kod testowy.

\subsection{Dane testowe syntetyczne}
TODO


\subsection{TPC}
TODO

\section{Porównanie sposobu formułowania kodu}

\end{document}
